\documentclass[a4paper,12pt,twoside]{memoir}

% Castellano
\usepackage[spanish,es-tabla]{babel}
\selectlanguage{spanish}
\usepackage[utf8]{inputenc}
\usepackage[T1]{fontenc}
\usepackage{lmodern} % Scalable font
\usepackage{microtype}
\usepackage{placeins}

\RequirePackage{booktabs}
\RequirePackage[table]{xcolor}
\RequirePackage{xtab}
\RequirePackage{multirow}

% Links
\usepackage[colorlinks]{hyperref}
\hypersetup{
	allcolors = {red}
}

% Ecuaciones
\usepackage{amsmath}

% Rutas de fichero / paquete
\newcommand{\ruta}[1]{{\sffamily #1}}

% Párrafos
\nonzeroparskip


% Imagenes
\usepackage{graphicx}
\newcommand{\imagen}[2]{
	\begin{figure}[!h]
		\centering
		\includegraphics[width=0.9\textwidth]{#1}
		\caption{#2}\label{fig:#1}
	\end{figure}
	\FloatBarrier
}

\newcommand{\imagenflotante}[2]{
	\begin{figure}%[!h]
		\centering
		\includegraphics[width=0.9\textwidth]{#1}
		\caption{#2}\label{fig:#1}
	\end{figure}
}



% El comando \figura nos permite insertar figuras comodamente, y utilizando
% siempre el mismo formato. Los parametros son:
% 1 -> Porcentaje del ancho de página que ocupará la figura (de 0 a 1)
% 2 --> Fichero de la imagen
% 3 --> Texto a pie de imagen
% 4 --> Etiqueta (label) para referencias
% 5 --> Opciones que queramos pasarle al \includegraphics
% 6 --> Opciones de posicionamiento a pasarle a \begin{figure}
\newcommand{\figuraConPosicion}[6]{%
  \setlength{\anchoFloat}{#1\textwidth}%
  \addtolength{\anchoFloat}{-4\fboxsep}%
  \setlength{\anchoFigura}{\anchoFloat}%
  \begin{figure}[#6]
    \begin{center}%
      \Ovalbox{%
        \begin{minipage}{\anchoFloat}%
          \begin{center}%
            \includegraphics[width=\anchoFigura,#5]{#2}%
            \caption{#3}%
            \label{#4}%
          \end{center}%
        \end{minipage}
      }%
    \end{center}%
  \end{figure}%
}

%
% Comando para incluir imágenes en formato apaisado (sin marco).
\newcommand{\figuraApaisadaSinMarco}[5]{%
  \begin{figure}%
    \begin{center}%
    \includegraphics[angle=90,height=#1\textheight,#5]{#2}%
    \caption{#3}%
    \label{#4}%
    \end{center}%
  \end{figure}%
}
% Para las tablas
\newcommand{\otoprule}{\midrule [\heavyrulewidth]}
%
% Nuevo comando para tablas pequeñas (menos de una página).
\newcommand{\tablaSmall}[5]{%
 \begin{table}
  \begin{center}
   \rowcolors {2}{gray!35}{}
   \begin{tabular}{#2}
    \toprule
    #4
    \otoprule
    #5
    \bottomrule
   \end{tabular}
   \caption{#1}
   \label{tabla:#3}
  \end{center}
 \end{table}
}

%
% Nuevo comando para tablas pequeñas (menos de una página).
\newcommand{\tablaSmallSinColores}[5]{%
 \begin{table}[H]
  \begin{center}
   \begin{tabular}{#2}
    \toprule
    #4
    \otoprule
    #5
    \bottomrule
   \end{tabular}
   \caption{#1}
   \label{tabla:#3}
  \end{center}
 \end{table}
}

\newcommand{\tablaApaisadaSmall}[5]{%
\begin{landscape}
  \begin{table}
   \begin{center}
    \rowcolors {2}{gray!35}{}
    \begin{tabular}{#2}
     \toprule
     #4
     \otoprule
     #5
     \bottomrule
    \end{tabular}
    \caption{#1}
    \label{tabla:#3}
   \end{center}
  \end{table}
\end{landscape}
}

%
% Nuevo comando para tablas grandes con cabecera y filas alternas coloreadas en gris.
\newcommand{\tabla}[6]{%
  \begin{center}
    \tablefirsthead{
      \toprule
      #5
      \otoprule
    }
    \tablehead{
      \multicolumn{#3}{l}{\small\sl continúa desde la página anterior}\\
      \toprule
      #5
      \otoprule
    }
    \tabletail{
      \hline
      \multicolumn{#3}{r}{\small\sl continúa en la página siguiente}\\
    }
    \tablelasttail{
      \hline
    }
    \bottomcaption{#1}
    \rowcolors {2}{gray!35}{}
    \begin{xtabular}{#2}
      #6
      \bottomrule
    \end{xtabular}
    \label{tabla:#4}
  \end{center}
}

%
% Nuevo comando para tablas grandes con cabecera.
\newcommand{\tablaSinColores}[6]{%
  \begin{center}
    \tablefirsthead{
      \toprule
      #5
      \otoprule
    }
    \tablehead{
      \multicolumn{#3}{l}{\small\sl continúa desde la página anterior}\\
      \toprule
      #5
      \otoprule
    }
    \tabletail{
      \hline
      \multicolumn{#3}{r}{\small\sl continúa en la página siguiente}\\
    }
    \tablelasttail{
      \hline
    }
    \bottomcaption{#1}
    \begin{xtabular}{#2}
      #6
      \bottomrule
    \end{xtabular}
    \label{tabla:#4}
  \end{center}
}

%
% Nuevo comando para tablas grandes sin cabecera.
\newcommand{\tablaSinCabecera}[5]{%
  \begin{center}
    \tablefirsthead{
      \toprule
    }
    \tablehead{
      \multicolumn{#3}{l}{\small\sl continúa desde la página anterior}\\
      \hline
    }
    \tabletail{
      \hline
      \multicolumn{#3}{r}{\small\sl continúa en la página siguiente}\\
    }
    \tablelasttail{
      \hline
    }
    \bottomcaption{#1}
  \begin{xtabular}{#2}
    #5
   \bottomrule
  \end{xtabular}
  \label{tabla:#4}
  \end{center}
}



\definecolor{cgoLight}{HTML}{EEEEEE}
\definecolor{cgoExtralight}{HTML}{FFFFFF}

%
% Nuevo comando para tablas grandes sin cabecera.
\newcommand{\tablaSinCabeceraConBandas}[5]{%
  \begin{center}
    \tablefirsthead{
      \toprule
    }
    \tablehead{
      \multicolumn{#3}{l}{\small\sl continúa desde la página anterior}\\
      \hline
    }
    \tabletail{
      \hline
      \multicolumn{#3}{r}{\small\sl continúa en la página siguiente}\\
    }
    \tablelasttail{
      \hline
    }
    \bottomcaption{#1}
    \rowcolors[]{1}{cgoExtralight}{cgoLight}

  \begin{xtabular}{#2}
    #5
   \bottomrule
  \end{xtabular}
  \label{tabla:#4}
  \end{center}
}


















\graphicspath{ {./img/} }

% Capítulos
\chapterstyle{bianchi}
\newcommand{\capitulo}[2]{
	\setcounter{chapter}{#1}
	\setcounter{section}{0}
	\chapter*{#2}
	\addcontentsline{toc}{chapter}{#2}
	\markboth{#2}{#2}
}

% Apéndices
\renewcommand{\appendixname}{Apéndice}
\renewcommand*\cftappendixname{\appendixname}

\newcommand{\apendice}[1]{
	%\renewcommand{\thechapter}{A}
	\chapter{#1}
}

\renewcommand*\cftappendixname{\appendixname\ }

% Formato de portada
\makeatletter
\usepackage{xcolor}
\newcommand{\tutor}[1]{\def\@tutor{#1}}
\newcommand{\course}[1]{\def\@course{#1}}
\definecolor{cpardoBox}{HTML}{E6E6FF}
\def\maketitle{
  \null
  \thispagestyle{empty}
  % Cabecera ----------------
\noindent\includegraphics[width=\textwidth]{cabecera}\vspace{1cm}%
  \vfill
  % Título proyecto y escudo informática ----------------
  \colorbox{cpardoBox}{%
    \begin{minipage}{.8\textwidth}
      \vspace{.5cm}\Large
      \begin{center}
      \textbf{TFG del Grado en Ingeniería Informática}\vspace{.6cm}\\
      \textbf{\LARGE\@title{}}
      \end{center}
      \vspace{.2cm}
    \end{minipage}

  }%
  \hfill\begin{minipage}{.20\textwidth}
    \includegraphics[width=\textwidth]{escudoInfor}
  \end{minipage}
  \vfill
  % Datos de alumno, curso y tutores ------------------
  \begin{center}%
  {%
    \noindent\LARGE
    Presentado por \@author{}\\ 
    en Universidad de Burgos --- \@date{}\\
    Tutor: \@tutor{}\\
  }%
  \end{center}%
  \null
  \cleardoublepage
  }
\makeatother

\newcommand{\nombre}{Raúl Negro Carpintero} %%% cambio de comando

% Datos de portada
\title{Sistema de Recomendación basado en Aprendizaje Profundo}
\author{\nombre}
\tutor{Bruno Baruque Zanón}
\date{\today}

\begin{document}

\maketitle


\newpage\null\thispagestyle{empty}\newpage


%%%%%%%%%%%%%%%%%%%%%%%%%%%%%%%%%%%%%%%%%%%%%%%%%%%%%%%%%%%%%%%%%%%%%%%%%%%%%%%%%%%%%%%%
\thispagestyle{empty}


\noindent\includegraphics[width=\textwidth]{cabecera}\vspace{1cm}

\noindent D. Bruno Baruque Zanón, profesor del departamento de nombre departamento, área de nombre área.

\noindent Expone:

\noindent Que el alumno D. \nombre, con DNI 71305764Z, ha realizado el Trabajo final de Grado en Ingeniería Informática titulado título de TFG. 

\noindent Y que dicho trabajo ha sido realizado por el alumno bajo la dirección del que suscribe, en virtud de lo cual se autoriza su presentación y defensa.

\begin{center} %\large
En Burgos, {\large \today}
\end{center}

\vfill\vfill\vfill

% Author and supervisor
\begin{minipage}{0.45\textwidth}
\begin{flushleft} %\large
Vº. Bº. del Tutor:\\[2cm]
D. nombre tutor
\end{flushleft}
\end{minipage}
\hfill
\begin{minipage}{0.45\textwidth}
\begin{flushleft} %\large
Vº. Bº. del co-tutor:\\[2cm]
D. nombre co-tutor
\end{flushleft}
\end{minipage}
\hfill

\vfill

% para casos con solo un tutor comentar lo anterior
% y descomentar lo siguiente
%Vº. Bº. del Tutor:\\[2cm]
%D. nombre tutor


\newpage\null\thispagestyle{empty}\newpage




\frontmatter

% Abstract en castellano
\renewcommand*\abstractname{Resumen}
\begin{abstract}
El tema principal son los sistemas de recomendación. El objetivo de este proyecto es comparar el rendimiento entre los modelos clásicos y los modelos basados en aprendizaje profundo sobre diferentes conjuntos de datos.
\end{abstract}

\renewcommand*\abstractname{Descriptores}
\begin{abstract}
Sistemas de recomendación, aprendizaje profundo, redes neuronales.
\end{abstract}

\clearpage

% Abstract en inglés
\renewcommand*\abstractname{Abstract}
\begin{abstract}
The main topic are the recommendation systems. The goal of this project is to compare the performance between classic models and those based on deep learning on different datasets.
\end{abstract}

\renewcommand*\abstractname{Keywords}
\begin{abstract}
Recommendation systems, deep learning, neural networks.
\end{abstract}

\clearpage

% Indices
\tableofcontents

\clearpage

\listoffigures

\clearpage

\listoftables
\clearpage

\mainmatter
\capitulo{1}{Introducción}
Los sistemas de recomendación son herramientas fundamentales de las cuales se aprovechan las principales compañías que ofrecen servicios para proveer a los usuarios de ítems relevantes.

Compañías como \textit{Amazon}, \textit{Netflix}, \textit{Spotify} y \textit{YouTube} utilizan la información de sus usuarios y de sus productos; así como las valoraciones que los propios usuarios hacen sobre los productos que consumen con el fin de crear modelos de recomendación capaces de ofrecer a estos usuarios ítems que les serán de interés la próxima vez que interactúen con el servicio.

Tal es la importancia de los sistemas de recomendación que, por ejemplo, \textit{Netflix}, \textit{Spotify} y \textit{Trivago} organizan competiciones en las que piden mejorar sus propios sistemas \cite{wiki:Netflix_Prize}, \cite{spotify}, \cite{trivago}. El premio que ofreció \textit{Netflix} por mejorar su modelo un 10\% fue de 1.000.000 de dólares.

Como se verá más adelante, existen varios tipos de sistemas de recomendación, siendo los más importantes:
\begin{itemize}
\tightlist
\item Modelos colaborativos
\item Modelos basados en contenido
\item Modelos híbridos
\end{itemize}
Los \textbf{modelos colaborativos} solo tienen en cuenta las valoraciones que los usuarios hacen sobre los ítems. Así, lo que hace en buscar usuarios parecidos (usuarios que valoran de forma parecida los mismos ítems) para recomendar ítems que uno no ha visto y el otro sí.

Los \textbf{modelos basados en contenido} tienen en cuenta las características de los ítems. De esta manera, buscan ítems con características parecidas a las de los ítems con los que ya ha interactuado el usuario positivamente. También hacen uso de las características de los usuarios.

Los \textbf{modelos híbridos} son una mezcla de los dos anteriores. Tienen en cuenta tanto las valoraciones como las características de usuarios e ítems.

En este proyecto serán objeto de estudio los modelos de recomendación clásico y los modelos de recomendación basados en aprendizaje profundo. Así pues, se evaluará cada modelo siguiendo unas métricas comunes a los dos para poder compararles.

\subsection{Estructura de la memoria}\label{estructura-memoria}
La memoria se divide en:
\begin{itemize}
\tightlist
\item \textbf{Introducción:} en este apartado se desarrolla de manera breve el tema que se va a tratar en el proyecto, así como la estructura del propio proyecto y los materiales entregados.
\item \textbf{Objetivos del proyecto:} sección en la que se indican los objetivos que se persiguen con la realización del proyecto, tanto técnicos como personales.
\item \textbf{Conceptos teóricos:} apartado en el que se explica todo lo necesario para el correcto entendimiento del tema tratado en el proyecto.
\item \textbf{Técnicas y herramientas:} sección en la que se listan todas las herramientas usadas en el proyecto, así como una breve justificación de su uso en favor de otras herramientas existentes.
\item \textbf{Aspectos relevantes del desarrollo del proyecto:} apartado en el que se explican temas de especial importancia.
\item \textbf{Trabajos relacionados:} sección en la que se indican y se desarrollan brevemente tanto artículos como proyectos que están directamente relacionados con este proyecto.
\item \textbf{Conclusiones y Líneas de trabajo futuras:} apartado en el que se recogen las conclusiones obtenidas una vez finalizado el proyecto, y las posibles mejoras que se pueden hacer en el futuro.
\end{itemize}

\subsection{Materiales adjuntos}\label{materiales-adjuntos}
Adicionalmente, junto a la memoria también se proporcionan los siguientes anexos:
\begin{itemize}
\tightlist
\item \textbf{Plan de Proyecto Software:} en este apéndice se desarrolla la planificación temporal que se ha seguido durante la realización del proyecto y la viabilidad del mismo.
\item \textbf{Especificación de Requisitos:} en este apéndice se indican y se explican los requisitos funcionales obtenidos a partir de los objetivos generales del proyecto.
\item \textbf{Especificación de diseño:} en este apéndice se explica cómo se ha estructurado el proyecto y cómo se han manipulado los datos de entrada para poder ser usados por los sistemas.
\item \textbf{Documentación técnica de programación:} esta sección contiene todo lo relacionado con la programación y la ejecución del proyecto.
\item \textbf{Documentación de usuario:} apéndice en el que se recoge todo lo que el usuario que va a utilizar la aplicación debe saber.
\end{itemize}
\capitulo{2}{Objetivos del proyecto}

\section{Objetivos técnicos}\label{objetivos-tecnicos}
\begin{itemize}
\tightlist
\item Comprender el funcionamiento de un sistema de recomendación clásico.
\item Recoger y evaluar los resultados obtenidos por el modelo clásico.
\item Comprender el funcionamiento de un sistema de recomendación basado en aprendizaje profundo.
\item Recoger y evaluar los resultados obtenidos por el modelo basado en aprendizaje profundo.
\item Comparar los resultados obtenidos por ambos modelos sobre una serie de conjuntos de datos.
\item Desplegar el programa Python en un servidor y hacer posible la interacción con nuevos usuarios.
\end{itemize}

\section{Objetivos personales}\label{objetivos-personales}
\begin{itemize}
\tightlist
\item Estudiar los sistemas de recomendación debido a su gran importancia.
\item Poner en uso los conocimientos adquiridos durante la carrera.
\item Ahondar más en el campo del aprendizaje profundo.
\end{itemize}



\capitulo{3}{Conceptos teóricos}

Los conceptos teóricos más importantes del proyecto son los relacionados con los sistemas de recomendación y las redes neuronales.

\section{Sistemas de recomendación}\label{sistemas-de-recomendacion}
Los sistemas de recomendación son herramientas que generan recomendaciones sobre un determinado objeto de estudio, a partir de las preferencias y opiniones dadas por los usuarios. \cite{sistemas-recomendacion}

En un sistema de recomendación hay dos clases de entidades: usuarios e ítems. Los datos están representados en una matriz de utilidad de tal manera que los usuarios son filas y los ítems columnas. Para cada par hay un valor que representa el grado de preferencia de un usuario para un ítem. Se asume que la mayoría de los valores se desconocen. Es por ello, que el objetivo de un sistema de recomendación es rellenar esos espacios en blanco.

\subsection{Tipos de sistemas de recomendación}\label{tipos-sistemas-recomendacion}
Existen tres tipos de sistemas (o modelos) de recomendación:
\begin{itemize}
\tightlist
\item Modelos basados en contenido
\item Modelos colaborativos
\item Modelos híbridos
\end{itemize}

\subsection{Modelos basados en contenido}\label{content-based-systems}
Este tipo de sistemas tienen en cuenta los gustos del usuario y las características de los ítems.

Para cada ítem hay que construir un perfil compuesto por las propiedades del mismo. Por ejemplo, si los ítem son películas, algunas de las propiedades serían género, director, fecha de estreno, reparto, etc. Una vez hecho esto, se recomiendan ítems cuyas propiedades sean parecidas a los ítems que el usuario ha valorado positivamente.

\subsection{Modelos colaborativos}\label{collaborative-filtering}
Recomiendan ítems basándose únicamente en los gustos de usuarios similares.

En este caso, en lugar de usar el vector ítem-perfil de un ítem, usamos las filas de la matriz de utilidad. Los usuarios son parecidos si sus filas se acercan de acuerdo a alguna distancia.

\subsection{Modelos híbridos}\label{modelos-hibridos}
Este tipo de sistemas utilizan las dos técnicas anteriores para ofrecer mejores recomendaciones.

Así pues, se recomendarán ítems que sean parecidos a los ítems valorados positivamente por usuarios parecidos.

\section{Medidas de calidad}\label{medidas-de-calidad}
Podemos hacer uso de diferentes métricas para conocer cómo de bueno es nuestro sistema de recomendación. Dado que las dos librerías que se han utilizado en el proyecto las ha creado la misma persona, se utilizarán las métricas ofrecidas en ellas por su gran similitud. Las métricas son:
\begin{itemize}
\tightlist
\item Precisión \textit{k}
\item Recall \textit{k}
\item Score \textit{AUC}
\item Ranking recíproco
\item \textit{Root Mean Square Error}
\end{itemize}

\subsection{Precisión \textit{k}}\label{precision-k}
La precisión \textit{k} \cite{precision_at_k} mide el número de elementos relevantes conocidos que se encuentran en las primeras \textit{k} posiciones del ranking de predicciones, es decir, el porcentaje de coincidencias entre los elementos relevantes conocidos y los elementos devueltos por el modelo. Su fórmula es:
\begin{equation}
Precision\;at\;k = \frac{ERC \subset kRP}{k}
\end{equation}
siendo \textit{ERC} los elementos relevantes conocidos y \textit{kRP} las primeras \textit{k} posiciones del ranking de predicciones.

\subsection{Recall \textit{k}}\label{recall-k}
El recall \textit{k} \cite{recall_at_k} la división del número de elementos relevantes que hay en las primeras \textit{k} posiciones del ranking de predicciones entre el número de elementos relevantes conocidos en el conjunto de test. Su fórmula es:
\begin{equation}
Recall\;at\;k = \frac{ERk}{ERT}
\end{equation}
siendo \textit{ERk} los elementos relevantes en las primeras \textit{k} posiciones del ranking de resultados y \textit{ERT} los elementos relevantes del conjunto de test.

\subsection{\textit{AUC} Score}\label{auc-score}
El \textit{AUC} (Area Under the Curve, Área Bajo la Curva) score \cite{auc_score} es la probabilidad de que un elemento relevante escogido aleatoriamente tenga un score mayor que un elemento no relevante escogido aleatoriamente.

\subsection{Ranking recíproco}\label{ranking-reciproco}
El ranking recíproco \cite{reciprocal_rank} es el inverso del score del elemento más relevante devuelto en el ranking de resultados. Su fórmula es:
\begin{equation}
Reciprocal\;rank = \frac{1}{SER}
\end{equation}
siendo \textit{SER} el score del elemento más relevante devuelto en el ranking de resultados.

En el modelo de \textit{Spotlight} aparece como \textit{MMR}, o \textit{Mean Reciprocal Rank}.

\subsection{Root Mean Square Error}\label{rmse}
El \textit{RMSE} se define como el error que hay entre dos conjuntos de datos. Aplicado a nuestro caso, el \textit{RMSE} compara un valor conocido con un valor observado. Su fórmula sería la siguiente:
\begin{equation}
RMSE = \sqrt{\frac{\sum_{i=1}^n(C_{i} - O_{i})^{2}}{n}}
\end{equation}
siendo \textit{C} una recomendación conocida y \textit{O} una recomendación aportada por el sistema.

\section{Tratamiento de los datos}\label{tratamiento-datos}
Lo más común en los sistemas de recomendación es tener una gran cantidad de usuarios y de ítems. En cambio, la cantidad de valoraciones de ítems por cada usuario que se tiene es muy pequeña. Esto hace que la matriz con la que representamos las valoraciones sea muy grande pero con muchos valores a 0 (o celdas vacías, dependiendo de la representación escogida).  

Esto es lo que llamamos \textit{matrices dispersas} \cite{wiki:Sparse_matrix}. Trabajar con estas matrices supone un gran coste de rendimiento, por lo que necesitamos convertirlas a otras estructuras con las que poder trabajar mejor.

En el caso de \textit{LightFM}, las estructuras utilizadas son:
\begin{itemize}
\tightlist
\item Matrices \textit{COO}
\item Matrices \textit{CSR}  
\end{itemize}

\subsection{Matrices COO}\label{matrices-coo}
Las \textit{matrices COO} \cite{coo-matrix}, o matrices coordenadas, dividen la matriz dispersa en tres vectores:
\begin{itemize}
\tightlist
\item Vector con los valores no nulos
\item Vector con el índice de las filas
\item Vector con el índice de las columnas
\end{itemize}
Esto lo podemos ver más claro con el siguiente ejemplo. Dada la siguiente matriz dispersa
\[\begin{bmatrix}
27&0&8&14&0&0\\
0&7&0&0&0&0\\
3&0&7&0&0&9\\
0&0&0&0&0&1\\
4&0&77&0&0&0\\
\end{bmatrix}\]
obtendríamos los siguientes vectores:
\begin{itemize}
\tightlist
\item Vector de valores: \[\begin{bmatrix} 27&8&14&7&3&7&9&1&4&77\\\end{bmatrix}\]
\item Vector de filas: \[\begin{bmatrix} 1&1&1&2&3&3&3&4&5&5\\\end{bmatrix}\]
\item Vector de columnas: \[\begin{bmatrix} 1&3&4&2&1&3&6&6&1&3\\
\end{bmatrix}\]
\end{itemize}

Este tipo de estructura es utilizada por \textit{LightFM} para representar la matriz de interacciones entre usuarios e ítems.

\subsection{Matrices CSR}\label{matrices-csr}
Otras estructuras ampliamente utilizadas, por ser más compactas que las matrices \textit{COO}, son las matrices \textit{CSR} \cite{csr-matrix}. Las matrices \textit{CSR} (\textit{Compressed Sparse Row}) no almacenan los índices de las filas, si no que almacenan punteros a los inicios de fila.

Utilizando como ejemplo la matriz dispersa del caso anterior, ahora obtendríamos:
\begin{itemize}
\tightlist
\item Vector de valores: \[\begin{bmatrix} 27&8&14&7&3&7&9&1&4&77\\\end{bmatrix}\]
\item Vector de columnas: \[\begin{bmatrix} 1&3&4&2&1&3&6&6&1&3\\
\end{bmatrix}\]
\item Vector de filas: \[\begin{bmatrix} 1&4&5&9&10&11\\\end{bmatrix}\]
\end{itemize}

Este tipo de estructura es utilizada por \textit{LightFM} para representar las matrices de features de usuarios e ítems.

\subsection{Spotlight}\label{spotlight}
\textit{Spotlight} utiliza directamente los arrays de \textit{numpy}, aunque también ofrece métodos para transformar estos arrays a matrices \textit{COO} y \textit{CSR}.


\capitulo{4}{Técnicas y herramientas}

\section{Metodologías}\label{metodologias}
\subsection{Scrum}\label{scrum}
Descripción

\section{Control de versiones}\label{control-de-versiones}
\subsection{Git}\label{git}
Descripción

\section{Repositorio}\label{repositorio}
\subsection{GitHub}\label{github}
Descripción

\section{Gestión del proyecto}\label{gestion-del-proyecto}
\subsection{ZenHub}\label{zenhub}
Descripción

\section{Entorno de Desarrollo Integrado (IDE)}\label{ide}
\subsection{Python}\label{python}
\begin{itemize}
\tightlist
\item Herramientas consideradas: 
	\href{https://jupyter.org/}{Jupyter}, 
	\href{https://www.spyder-ide.org/}{Spyder} y
	\href{https://www.jetbrains.com/pycharm/}{PyCharm}
\item Herramienta escogida:
	Jupyter \href{https://jupyter.org/}{Jupyter}
\end{itemize}
Justificación

\subsection{Documentación}\label{documentacion}
\begin{itemize}
\tightlist
\item Herramientas consideradas: 
	\href{http://www.xm1math.net/texmaker/}{Texmaker} y
	\href{https://www.openoffice.org/es/}{Apache OpenOffice}
\item Herramienta escogida:
	\href{http://www.xm1math.net/texmaker/}{Texmaker}
\end{itemize}
Justificación

\section{Librerías}\label{librerias}
\subsection{NumPy}\label{numpy}
\href{http://www.numpy.org/}{NumPy} descripción

\subsection{SciPy}\label{scipy}
\href{https://www.scipy.org/scipylib/index.html}{SciPy} descripción

\subsection{LightFM}\label{lightfm}
\href{https://github.com/lyst/lightfm}{LightFM} descripción

\section{Datasets}\label{datasets}
\subsection{Anime}\label{anime}
Este conjunto de datos \cite{CopperUnion2017} contiene información relativa a las preferencias de 73.516 usuarios sobre 12.294 películas o series de anime. Cada usuario ha valorado unos 106 animes.

El conjunto de datos se encuentra disponible en: \href{https://www.kaggle.com/CooperUnion/anime-recommendations-database/downloads/anime-recommendations-database.zip/1}{Anime Dataset}.

\subsection{Book-Crossing}\label{book-crossing}
Este conjunto de datos \cite{Ziegler2004} contiene 1.149.780 valoraciones de 278.858 usuarios sobre 271.379 libros. Hay alrededor de 11 valoraciones por usuario.

El conjunto de datos se encuentra disponible en \href{http://www2.informatik.uni-freiburg.de/~cziegler/BX/BX-CSV-Dump.zip}{Book-Crossing Dataset}.

\subsection{MovieLens}\label{movielens}
Este conjunto de datos \cite{GroupLens1998} contiene 100.000 valoraciones de 943 usuarios sobre 1.682 películas. Hay unas 106 valoraciones por usuario.

El conjunto de datos se encuentra disponible en los propios archivos descargados por la librería LightFM \cite{DBLP:conf/recsys/Kula15}.

\subsection{Last.FM}\label{last.fm}
Este conjunto de datos \cite{Cantador:RecSys2011} contiene las veces que 1.892 usuarios han escuchado a 17.632 artistas. Cada usuario ha escuchado alrededor de unos 49 artistas.

El conjunto de datos se encuentra disponible en \href{http://files.grouplens.org/datasets/hetrec2011/hetrec2011-lastfm-2k.zip}{Last.FM}.

\subsection{Dating Agency}\label{dating-agency}
Este conjunto de datos \cite{brozovsky07recommender} contiene 17.359.346 valoraciones de 135.359 usuarios sobre 168.791 perfiles de otros usuarios. Hay unas 128 valoraciones por usuario.

El conjunto de datos se encuentra disponible en \href{http://www.occamslab.com/petricek/data/libimseti-complete.zip}{Dating Agency}.
\capitulo{5}{Aspectos relevantes del desarrollo del proyecto}

Este apartado pretende recoger los aspectos más interesantes del desarrollo del proyecto.

\section{Metodologías}\label{metodologias}
Para la gestión del proyecto se ha adaptado la metodología ágil Scrum. Las principales características han sido:
 \begin{itemize}
\tightlist
\item
	La duración de los sprints fue de dos semanas aproximadamente.
\item
	Los sprints finalizaban con una reunión en la que se revisaba el trabajo hecho y se planteaban las tareas del siguiente sprint.
\item	
	Se utilizó ZenHub para llevar el seguimiento de las tareas.
\end{itemize}

\section{Formación}\label{formacion}
El proyecto ha requerido obtener una serie de conocimientos que no se tenían inicialmente. Se ha estudiado el campo de los sistemas de recomendación y la utilización de aprendizaje profundo en Python.

Para la formación en los sistemas de recomendación se usó el libro:
\begin{itemize}
\tightlist
\item
	\emph{Mining of Massive Datasets} (Jure Leskovec, Anand Rajaraman y Jeffrey D. Ullman) \cite{miningDatasets}.
\end{itemize}

Para la formación en aprendizaje profundo en Python se realizó el siguiente curso:
\begin{itemize}
\tightlist
\item
	\emph{Practical Deep Learning For Coders, v3} (Jeremy Howard) \cite{fastai}.
\end{itemize}

Para poder usar la aplicación desde una interfaz web se siguió la documentación de \textit{Flask} y de \textit{WTForms}:
\begin{itemize}
\tightlist
\item
	\emph{Welcome to Flask} (Pallets Team) \cite{flask}.
\item
	\emph{WTForms Documentation} (WTForms Team) \cite{wtforms}.
\end{itemize}

Para la obtención de los modelos de recomendación se utilizaron las siguientes librerías:
\begin{itemize}
\tightlist
\item
	\emph{LightFM} (Maciej Kula) \cite{kulalightfm}.
\item
	\emph{Spotlight} (Maciej Kula) \cite{kula2017spotlight}.
\end{itemize}

\section{Desarrollo del código}\label{desarrollo-codigo}
El proyecto se ha centrado en dos aspectos fundamentales: obtener un modelo clásico y obtener un modelo basado en aprendizaje profundo.

Para la obtención del modelo clásico se estudió en primer lugar hacer uso de la librería \textit{Crab}. Esta opción se descartó debido a la cantidad de fallos que se dieron durante la instalación (no se pudo llegar a instalar) y a la falta de actividad por parte de sus autores en el repositorio (no se ha tocado desde hace 7 años). Por todo esto, se decidió utilizar la librería \textit{LightFM}.

Para la obtención del modelo basado en aprendizaje profundo se estudió en primer lugar hacer uso de la librería de \textit{fast.ai} \cite{fastai}, pero se descartó porque \textit{Spotlight} parecía ser más amigable en cuanto a facilidad de uso; y porque está creado por la misma gente que creó \textit{LightFM}, lo cual hace que la manera de trabajar con estas dos herramientas sea muy parecida. Así pues, se espera que la evaluación de los modelos siga los mismos criterios en ambas librerías.

Para la obtención de la aplicación web se utilizó \textit{Flask}. Es una librería que aparece en una gran cantidad de artículos elogiando su facilidad de uso y su potencial. Cumplió todas las funciones que se esperaba de ella. Junto con \textit{Flask} se utilizó \textit{WTForms} para crear todos los formularios que forman parte de la aplicación.

Todo el código del proyecto se pensó para seguir el \textit{Modelo Vista Controlador}.

\section{CUDA}\label{cuda}
Para la obtención y el entrenamiento de los modelos de \textit{Spotlight} es posible utilizar la GPU gracias a \textit{CUDA}. Si el equipo en el que corre la aplicación no tuviera GPU, o no estuviera disponible, no pasaría nada (no saltarían errores), simplemente se ejecutarían las operaciones en la CPU.

La instalación de \textit{Spotlight} incluye \textit{CUDA}.

\section{Problemas}\label{problemas}
Durante el desarrollo del proyecto se han encontrado los siguientes problemas:

\subsection{Predicciones en modelo de secuencia}\label{predicciones-secuencia}
De momento, no se ha conseguido obtener las predicciones para el modelo de secuencia de \textit{Spotlight}. Esto es debido a que a diferencia del resto de modelos, con los cuales para ver las predicciones basta con indicar el usuario cuya predicciones se quiere obtener, en el modelo de secuencia no vale el usuario. Hay que utilizar una secuencia de interacciones del usuario. Esto dificulta el trabajo a la hora de programarlo, además de complicarle la vida al usuario que va a interactuar con la aplicación.

En lugar de que la persona que vaya a probar la aplicación tenga que elegir el id del usuario cuyas predicciones quiere ver, tendría que saber de ante mano la secuencia de interacciones que ese usuario ha tenido con los ítems para que el modelo pueda calcular los siguientes con los que podría interactuar.

\subsection{Añadir valoraciones}\label{mas-valoraciones}
Tampoco ha sido posible llevar a cabo la adición de más valoraciones por parte de usuarios ya existentes o nuevos. No queda muy claro cómo proceder con los métodos de \textit{fit\_partial} de los modelos. Además, según han comentado algunos usuarios de las librerías, parece que hay que tener en cuenta con antelación si se va a querer añadir nuevos usuarios, porque deberían existir en el conjunto de datos original (aunque no interactúen con nada) \cite{problema-mas-valoraciones}.

\subsection{Heroku}\label{heroku}
Se ha intentado desplegar la aplicación \textit{Flask} a través de \textit{Heroku} y \textit{Gunicorn}, pero no ha sido posible debido a que la aplicación de \textit{Heroku} no era capaz de trabajar con los paquetes del proyecto. Siempre daba fallo de \textit{ModuleNameNotFound}.

Se ha intentado cambiar la forma de importar los paquetes sin suerte.

\section{Documentación}\label{documentacion}
Las opciones que se plantearon para realizar la documentación fueron \textit{Apache OpenOffice} y \textit{Texmaker}. Ya que el Trabajo Fin de Grado se puede ver como una forma de aprender conceptos nuevos que se desmarcan de lo visto durante la carrera, se optó por \textit{Texmaker} debido a la novedad y al querer aprender y usar \LaTeX{}.

Además, la documentación se apoya en el programa \textit{JabRef} \cite{jabref}, con el que se va guardando un registro con las entradas \textit{bibtex} de los artículos y demás elementos que se han consultado a lo largo del proyecto.

\capitulo{6}{Trabajos relacionados}

Este apartado sería parecido a un estado del arte de una tesis o tesina. En un trabajo final grado no parece obligada su presencia, aunque se puede dejar a juicio del tutor el incluir un pequeño resumen comentado de los trabajos y proyectos ya realizados en el campo del proyecto en curso. 

\capitulo{7}{Conclusiones y Líneas de trabajo futuras}

\section{Conclusiones}\label{conclusiones}
Como se ha podido observar a lo largo de todos los artículos que se han leído, y por la experiencia personal con determinados servicios, está claro que los sistemas de recomendación son unas herramientas muy potentes y que tienen una importancia mayúscula en nuestras vidas. Hemos normalizado completamente que compañías como \textit{Neflix}, \textit{Amazon}, \textit{Spotify}, etc. nos ofrezcan productos o servicios que piensan que nos pueden gustar y que, de hecho, en muchos casos nos acaban gustando.

A continuación se muestran parte de los resultados obtenidos por los distintos modelos con el conjunto de datos de \textit{Movielens}, ya que al haber realizado cambios de última hora en los sistemas, es inviable entrenar a tiempo todos los modelos varias veces con todos los datasets de prueba:

\tablaSinColores{Métricas modelos \textit{LightFM}}{l l l l l}{5}{lightfm}
{\textbf{Modelo} & \textbf{Precisión k} & \textbf{AUC Score} & \textbf{Recall k} & \textbf{Ranking recíproco} \\}{
	Colaborativo & 0,1879 & 0,8675 & 0,1218 & 0,4304 \\
	Híbrido & 0,1759 & 0,8655 & 0,1160 & 0,4017 \\
	Por contenido & 0,1746 & 0,8640 & 0,1143 & 0,3910 \\
}

Como se puede observar, el modelo colaborativo es el que mejor sale evaluado. Esto contrasta con la teoría, ya que el modelo híbrido debería ser mejor por utilizar tanto las valoraciones como las features de los usuarios y de los ítems.

Se tendría que repetir las pruebas con otros conjuntos de datos distintos para ver si la tendencia es esa o cambia.

\tablaSinColores{Métricas modelos \textit{Spotlight}}{l l l l l}{5}{spotlight}
{\textbf{Modelo} & \textbf{RMSE} & \textbf{MRR} & \textbf{Precisión k} & \textbf{Recall k} \\}{
	Fact. expl. time & 1,1457 & 0,0102 & 0,0424 & 0,0168 \\
	Fact. expl. & 1,1471 & 0,0114 & 0,0409 & 0,0177 \\
	Fact. impl. time & - & 0,0625 & 0,2317 & 0,1380 \\
	Fact. impl. & - & 0,0671 & 0,2367 & 0,1482 \\
	Secuencia impl. - & 0,0304 & - & - \\
}

Como se puede observar, los mejores modelos son los que no hacen uso de los timestamps. Entre ellos, el que mejor sale parado es el modelo de factorización implícito sin timestamps (tampoco hace uso de las valoraciones).

Si comparamos ambos modelos:

\tablaSinColores{\textit{LightFM} vs \textit{Spotlight}}{l l l l}{4}{versus}
{\textbf{Modelo} & \textbf{Precisión k} & \textbf{Recall k} & \textbf{MRR} \\}{
	Colaborativo & 0,1879 & 0,1218 & 0,4304 \\
	Fact. impl. & 0,2367 & 0,1482 & 0,0671 \\
}

Como se puede observar, el modelo de factorización implícito obtiene mejores valores de precisión y recall que el modelo colaborativo, pero sale perdiendo, y por mucho, en el ranking recíproco.

Con estos resultados, se lleva a la conclusión de que el mejor modelo de \textit{LightFM} es el que solo utiliza las valoraciones para recomendar, mientras que el mejor modelo de \textit{Spotlight} es uno de los que no utilizan las valoraciones para nada.

A parte de estos resultados evidentes, se ha podido comparar el tiempo que tarda uno de los sistemas de \textit{LightFM} en leer el conjunto de \textit{Dating Agency}, obtener, entrenar y evaluar el modelo frente a lo que tarda un modelo de \textit{Spotlight} en hacer lo mismo. Mientras que el modelo clásico tardaba unas 12 horas en completar las tareas, el modelo basado en aprendizaje profundo lo ha hecho en unas 4 horas. Esto es lo que podemos esperar de los modelos de deep learning gracias al uso de \textit{CUDA}.

\section{Líneas de trabajo futuras}\label{lineas-futuras}
Dado que no se han podido obtener las predicciones para el modelo de secuencia, se podría plantear como futura línea de trabajo obtenerlas.

También se deja como futura mejora el añadir valoraciones de usuarios existentes o de usuarios nuevos.

Desgraciadamente, ni \textit{LightFM} ni \textit{Spotlight} implementan métodos de búsqueda de hiperparámetros. Tampoco lo hace \textit{PyTorch}, por lo que podría ser interesante el crear un método capaz de obtener los mejores parámetros para los modelos, ya sea una búsqueda aleatoria o en malla.

Otra posible tarea es la de desplegar la aplicación en un servidor con hardware "competitivo" y comprar los resultados para determinar hasta qué punto el hardware del equipo ralentiza la obtención y el entrenamiento de los modelos. La mejora va a ser muy notable sí o sí. Tan solo hay que ver que los modelos clásicos tardan hasta 12 horas en completar todo el proceso ejecutándose en la CPU (Intel Core i5-6000) mientras que los basados en aprendizaje profundo tardan unas x horas ejecutándose en una GPU de gama media-alta pero con ya unos años a sus espaldas(nVidia GTX 970). Es por eso que sería interesante ver el desempeño de los modelos basados en deep learning en hardware más moderno y potente (tal ver una nVidia GTX Titan o Titan V).

Dado que se han realizado cambios de última hora en la obtención de los sistemas, estaría bien en un futuro el poder tener una comparativa más exhaustiva y extensa de las métricas de cada modelo.

También se deja como futura mejora la creación de tests y pruebas sobre el código, que no ha sido posible realizar debido a la falta de tiempo.


\bibliographystyle{plain}
\bibliography{bibliografia}

\end{document}
